3. (2 pts) Sea \( V \) un espacio vectorial sobre el campo \( F \) y considere \( f: V \rightarrow V \) una transformación lineal. Demuestre que son equivalentes: \\
a) \( \operatorname{Im}(f)=\operatorname{Im}\left(f^{2}\right)\left(\right. \) donde \( \left.f^{2}=f \circ f\right) \). \\
b) \( V=N u c(f) \oplus \operatorname{Im}(f) \). \\

\begin{solution} \\

$\Rightarrow )$\\
Supongamos que \( \operatorname{Im}(f) = \operatorname{Im}(f^2) \).


Definamos f: Consideremos \( f: V \rightarrow V \) tal que para cualquier \( v \in V \), \( f(v) \in \operatorname{Im}(f) \).


Propiedad de $Im(f^2)$: Dado que \( \operatorname{Im}(f) = \operatorname{Im}(f^2) \), para cualquier \( w \in \operatorname{Im}(f) \), existe \( v \in V \) tal que \( f(v) = w \) y \( f^2(v) = f(w) \).


Ahora, consideremos cualquier \( v \in V \). Tenemos:
\[
   v = (v - f(v)) + f(v)
   \]
Aquí, \( f(v) \in \operatorname{Im}(f) \) y \( v - f(v) \in \operatorname{Nuc}(f) \) porque \( f(v - f(v)) = f(v) - f(v) = 0 \), este paso es dificil de ver pero recordemos que $f$ es una transformación lineal así que cumple superposición, además recordemos que la imagen cuadrada de la función es igual a la imagen de la función por ello que lleguemos a que esta diferencia es igual a cero.


Para establecer que \( V = \operatorname{Nuc}(f) \oplus \operatorname{Im}(f) \), debemos demostrar que la intersección de \( \operatorname{Nuc}(f) \) y \( \operatorname{Im}(f) \) es solo el vector cero. Si \( x \in \operatorname{Nuc}(f) \cap \operatorname{Im}(f) \), entonces \( f(x) = 0 \) y \( x = f(y) \) para algún \( y \in V \). Pero esto implica que \( f(f(y)) = f(x) = 0 \), lo que significa que \( f(y) \in \operatorname{Nuc}(f) \), y por lo tanto \( x = 0 \).


\( \therefore V = \operatorname{Nuc}(f) \oplus \operatorname{Im}(f) \)
\quad\quad\quad\quad\quad\quad\quad\quad\quad\quad\quad\quad\quad\quad\quad\quad\quad\quad\quad\quad\quad\quad\quad\quad\quad\quad\quad\quad\quad\quad\quad\quad\quad\quad\quad \triangledown \\

$\Leftarrow )$ \\
Supóngase que \( V = \operatorname{Nuc}(f) \oplus \operatorname{Im}(f) \)


Cualquier vector \( v \in V \) puede escribirse de manera única como \( v = u + w \) con \( u \in \operatorname{Nuc}(f) \) y \( w \in \operatorname{Im}(f) \). Dado que \( f(u) = 0 \) y \( f(w) \) se mantiene en \( \operatorname{Im}(f) \), es claro que cualquier imagen bajo \( f \) de un elemento en \( \operatorname{Im}(f) \) también está en \( \operatorname{Im}(f) \). \\
Dado que \( \operatorname{Im}(f) \subseteq \operatorname{Im}(f^2) \) por definición, y \( \operatorname{Im}(f^2) \subseteq \operatorname{Im}(f) \) por hipótesis de la suma directa, tenemos \( \operatorname{Im}(f) = \operatorname{Im}(f^2) \). \\

Así, demostramos que si \( V = \operatorname{Nuc}(f) \oplus \operatorname{Im}(f) \), entonces \( \operatorname{Im}(f) = \operatorname{Im}(f^2) \). \end{solution}